\section{Introduction}

Sign-language translation requires extracting linguistic evidence from motion and
turning it into text.  LLMs can address the latter, but visual encoders need a
well-defined interface.  The first published \system{} mapped 2D keypoints to a
frozen LLM's token-embedding space and reported preliminary MSE--cosine loss near
$8\times10^{-4}$ \citep{velasquez2026imitator}.  Yet geometric proximity does not
guarantee the intended token, particularly in anisotropic contextual spaces
\citep{ethayarajh2019contextual}, and its splits were not signer-exclusive.

We preserve the premise---a compact model learns the input language of a frozen
LLM---but predict exact \Gemma{} tokenizer IDs.  Outputs are directly measurable
and can enter the ordinary \Gemma{} 3n text interface \citep{gemma3n} for
downstream linguistic realization.

We first ask whether the visual model recognizes one isolated sign and emits
its exact \Gemma{} IDs.  As a step toward the original goal of sentence
translation, we then ask how much of that ability carries over to ordered
transcription of two to eight signs.  Randomly concatenating clips lets us test
visual ordering without a linguistic next-sign prior.  The sequence benchmark
uses signer-disjoint folds, two prospectively held-out signers, and segment
permutation as a control for visual order.  A separate exploratory audit asks
whether \Gemma{} can recover a corrupted label without sentence context.
